\documentclass[11pt]{article}
\usepackage[margin=1in]{geometry}
\usepackage{amsmath,amssymb,amsthm}
\usepackage{graphicx}
\usepackage{booktabs}
\usepackage{hyperref}
\usepackage{natbib}
\usepackage{xcolor}
\usepackage{microtype}
\usepackage{algorithm}
\usepackage{algpseudocode}

\newcommand{\Gr}{\mathrm{Gr}}
\newcommand{\BN}{\mathrm{BN}}
\newcommand{\R}{\mathbb{R}}
\newcommand{\TODO}[1]{\textcolor{red}{[TODO: #1]}}
\newcommand{\CITE}[1]{\textcolor{red}{[CITE: #1]}}

\title{Causal Geometry Audit of RNA Foundation Model Representations:\\
  Bracket Norm, Grassmannian Transportability, and Weight Factorization\\
  on HTT mRNA for Huntington's Disease}

\author{Elliot Tower\\
  Mechanistic Research\\
  \texttt{elliot@elliottower.ai}}

\date{July 2026}

\begin{document}
\maketitle

\begin{abstract}
RNA-targeting therapeutics face two geometric confounds: distinguishing
genuine structural features from length-dependent artifacts, and
predicting whether drug mechanisms transfer across pathogenic mutations.
We apply three independent audits to RNA-FM, a 99M-parameter RNA
foundation model, using Huntingtin (HTT) mRNA as a case study.
%
First, bracket-norm correction reveals which embedding dimensions encode
intensive structural properties of the CAG repeat region versus
extensive properties that scale with sequence length.
%
Second, Grassmannian subspace analysis quantifies how the
representation subspace drifts as CAG repeat count increases from
wild-type (21 repeats) through the disease threshold (36+) to severe
expansions (80 repeats), yielding a transportability score for
antisense oligonucleotide design.
%
Third, singular value decomposition of RNA-FM's weight matrices reveals
progressive compression: $W_Q$ effective rank drops from 265 (layer~0)
to 86 (layer~11), while $W_Q$--$W_K$ subspace overlap reaches 0.81 in
layer~3, confirming that the query-key circuit reads from a shared
low-dimensional subspace amenable to mechanistic decomposition.
%
All analyses, code, and data are publicly available.
\end{abstract}

%% ────────────────────────────────────────────
\section{Introduction}
\label{sec:intro}

Huntington's disease (HD) is caused by expansion of a CAG trinucleotide
repeat in exon~1 of the Huntingtin gene (HTT), producing a
polyglutamine-expanded protein that aggregates and causes
neurodegeneration~\citep{macdonald1993}.
Normal alleles carry 6--35 CAG repeats; 36--39 repeats confer incomplete
penetrance; 40+ repeats are fully penetrant~\citep{rubinsztein1996}.
No disease-modifying therapy exists. The primary therapeutic strategy
targets HTT mRNA directly, using antisense oligonucleotides (ASOs) to
reduce production of the toxic protein~\citep{tabrizi2019}.

The central challenge for RNA-targeting drug design is \emph{transportability}:
a drug mechanism validated on one CAG repeat length must transfer to the
repeat lengths seen in patients. Clinical trials of tominersen
(an ASO targeting HTT mRNA) were halted partly because the drug reduced
both mutant and wild-type HTT with insufficient selectivity~\citep{tabrizi2022}.
Understanding which structural features of the CAG repeat region are
preserved across repeat lengths---and which degrade---would inform
the design of allele-selective therapies.

RNA foundation models offer a new lens on this problem. RNA-FM, a
99M-parameter BERT-style model pre-trained on 23 million non-coding RNA
sequences, produces 640-dimensional contextual embeddings for each
nucleotide position~\citep{rnafm}. These embeddings
capture structural and functional properties learned from evolutionary
data, providing a representation space where geometric analyses can
detect genuine structural signals.

We apply three complementary audits from the causal geometry
framework~\citep{tower2026bn,tower2026grass}
to RNA-FM embeddings of the HTT CAG repeat region:

\begin{enumerate}
\item \textbf{Bracket-norm correction.} RNA embeddings are confounded
  by sequence length: longer sequences produce higher-norm
  representations due to positional encoding and attention aggregation
  effects. The bracket norm~\citep{tower2026bn} normalizes
  feature magnitudes by $\sqrt{n}$ to distinguish \emph{intensive}
  structural properties (genuine features of the RNA fold) from
  \emph{extensive} properties (artifacts scaling with input length).

\item \textbf{Grassmannian transportability.} For each CAG repeat count
  $r$, we compute the principal component subspace of RNA-FM embeddings
  on the repeat region, yielding a point on the Grassmannian
  $\Gr(k, 640)$. The geodesic distance between the wild-type subspace
  ($r = 21$) and a mutant subspace ($r > 36$) quantifies how much the
  representation structure degrades---a direct measure of drug
  mechanism transferability~\citep{tower2026grass}.

\item \textbf{Weight factorization.} Singular value decomposition of
  RNA-FM's attention and MLP weight matrices reveals the intrinsic
  dimensionality of each projection. Cross-projection subspace overlap
  (shared singular directions between $W_Q$ and $W_K$, between $W_V$
  and $W_O$) indicates whether the model's internal structure admits a
  shared factor bank decomposition, as demonstrated for GPT-2 in prior
  work~\citep{tower2026fac}.
\end{enumerate}

Together, these three audits provide a geometric validity check on any
structural claim derived from RNA-FM embeddings: features surviving
bracket-norm correction, Grassmannian transportability, and weight-space
factorization analysis carry higher evidential weight than features
identified by any single method.


%% ────────────────────────────────────────────
\section{Background}
\label{sec:background}

\subsection{RNA-FM architecture}

RNA-FM~\citep{rnafm} is a BERT-style~\citep{devlin2019bert}
transformer encoder with 12 layers, 20 attention heads per layer,
hidden dimension $d = 640$, and intermediate MLP dimension $d_{\mathrm{ff}} = 5120$.
The model uses a nucleotide-level vocabulary of 28 tokens (A, C, G, U,
plus special tokens) and was pre-trained with masked language modeling
on 23.7 million non-coding RNA sequences from RNAcentral.
%
RNA-FM achieves state-of-the-art performance on RNA secondary structure
prediction, producing per-nucleotide embeddings that encode structural
context without explicit thermodynamic modeling.

\subsection{HTT mRNA and CAG repeat structure}

The human HTT mRNA (RefSeq NM\_002111.7) spans 13{,}669 nucleotides.
The pathogenic CAG repeat lies in exon~1, with normal alleles containing
approximately 21 repeats (63 nucleotides). The repeat region and its
flanking sequences create a structured RNA domain whose conformation
depends on repeat length: longer repeats form more stable hairpin
structures that may affect mRNA splicing, localization, and
translation~\citep{demezer2011,sobczak2003}.

\subsection{Bracket norm}

The bracket norm~\citep{tower2026bn} is a normalization that
distinguishes intensive from extensive properties in high-dimensional
feature spaces. For a feature vector $\mathbf{f} \in \R^d$ computed
from $n$ observations, the bracket-norm corrected value is
$\BN(\mathbf{f}) = \|\mathbf{f}\| / \sqrt{n}$.
%
Features whose bracket-norm corrected magnitude is invariant to $n$
represent genuine intensive properties; features whose corrected
magnitude varies with $n$ are extensive and may be confounded by
sampling density or input size. In the RNA context, $n$ is the
sequence length, and the correction identifies embedding dimensions
that encode per-nucleotide structural information rather than
global sequence statistics.

\subsection{Grassmannian subspace analysis}

A $k$-dimensional subspace of $\R^d$ is a point on the Grassmannian
manifold $\Gr(k, d)$~\citep{edelman1998}. Given two subspaces
$U, V \in \Gr(k, d)$, the principal angles
$\theta_1 \leq \cdots \leq \theta_k$ between them are computed via the
SVD of $U^\top V$, and the geodesic distance is
$d_g(U, V) = \sqrt{\sum_i \theta_i^2}$.
%
This distance quantifies structural similarity between representation
subspaces: small geodesic distance implies that the same geometric
features are active in both conditions, while large distance indicates
a qualitative shift in representation structure.
%
When applied to RNA embeddings across CAG repeat variants, this yields
a \emph{transportability score}: the degree to which drug-relevant
structural features identified on wild-type RNA persist in disease-length
expansions.


%% ────────────────────────────────────────────
\section{Methods}
\label{sec:methods}

\subsection{Data preparation}

We obtained the full-length HTT mRNA sequence (NM\_002111.7,
13{,}669~nt) from NCBI RefSeq. The native CAG repeat region was
identified by regex matching, yielding 21 tandem CAG repeats at
positions 148--211.

To study transportability across repeat lengths, we constructed
synthetic HTT mRNA variants by replacing the native CAG repeat with
$r \in \{15, 21, 27, 36, 40, 50, 60, 80\}$ tandem repeats, preserving
50~nt of flanking sequence on each side. This set spans
sub-normal (15), normal (21), intermediate (27), reduced-penetrance
threshold (36), full-penetrance (40), and severe expansion (50--80)
alleles.

\subsection{Embedding extraction}

We loaded RNA-FM weights from the HuggingFace
repository into a BERT encoder using the
\texttt{transformers} library~\citep{wolf2020transformers}, with
architecture parameters matched to the published
specification (12 layers, 640 hidden, 20 heads, 5120 intermediate).
%
Input sequences were tokenized at single-nucleotide resolution
(A$\to$5, C$\to$6, G$\to$7, U$\to$8) with CLS and EOS tokens.
Embeddings from the final hidden layer and all intermediate layers
were extracted for each variant.

\subsection{Bracket-norm analysis}

For windows of size $n \in \{100, 200, 300, 400, 500\}$~nt centered on
the CAG repeat region, we computed the mean embedding norm
$\bar{\ell}(n) = \frac{1}{n} \sum_{i=1}^{n} \|\mathbf{h}_i\|$
and the bracket-norm corrected value $\BN(n) = \bar{\ell}(n) / \sqrt{n}$.
%
An intensive property yields constant $\BN(n)$ across window sizes;
an extensive property yields $\BN(n)$ that decreases as $1/\sqrt{n}$.

We also computed per-dimension bracket norms: for each of the 640
embedding dimensions, we computed the $L^2$ norm across positions and
divided by $\sqrt{n}$, identifying the dimensions with the highest
intensive signal.

\subsection{Grassmannian transportability}

For each CAG repeat variant $r$, we computed RNA-FM embeddings on the
synthetic sequence and extracted the final-layer representations for
the repeat-plus-flanking region. PCA with $k = 10$ components yielded
a subspace $U_r \in \Gr(10, 640)$.

The geodesic distance $d_g(U_{21}, U_r)$ between the wild-type
subspace and each variant subspace quantifies transportability.
We also computed the full principal angle spectrum to identify which
subspace dimensions are preserved versus disrupted by repeat expansion.

\subsection{Layer-wise analysis}

For a 500~nt window centered on the CAG repeat, we extracted hidden
states from all 13 layers (embedding + 12 transformer layers) and
performed PCA on each. The effective dimensionality (number of
components for 95\% variance) characterizes how concentrated the
representation is at each processing stage.

We computed cross-layer subspace similarity as the mean cosine of
principal angles between 10-dimensional PCA subspaces of different
layers, yielding a $13 \times 13$ similarity matrix.

\subsection{Weight matrix factorization}

For each of the 12 transformer layers, we performed SVD on all six
projection matrices ($W_Q, W_K, W_V, W_O \in \R^{640 \times 640}$
and $W_{\mathrm{in}} \in \R^{5120 \times 640}$,
$W_{\mathrm{out}} \in \R^{640 \times 5120}$).
%
The effective rank at 95\% cumulative variance and the spectral decay
ratio $\sigma_1 / \sigma_{10}$ characterize the intrinsic
dimensionality of each projection.

To assess whether projections share a common factor structure, we
computed the subspace overlap between the top-20 right singular vectors
of each attention projection pair ($W_Q$--$W_K$, $W_Q$--$W_V$,
$W_K$--$W_V$, $W_Q$--$W_O$) within each layer. High overlap indicates
that the same input-space directions are important across projections,
supporting a shared factor bank decomposition.


%% ────────────────────────────────────────────
\section{Results}
\label{sec:results}

\subsection{LayerNorm absorbs per-position magnitude variation}

The mean per-position embedding norm is $28.667 \pm 0.001$ across all
window sizes (100--500~nt), yielding a bracket-norm corrected value
that decreases monotonically as $1/\sqrt{n}$ (from 2.867 at $n=100$
to 1.282 at $n=500$). This constancy reflects the LayerNorm applied
at every transformer layer: magnitude information is normalized away,
and structural differences between nucleotide positions are encoded
entirely in the \emph{directions} of the embedding vectors.

Per-dimension bracket norms identify a sparse set of intensive
dimensions. The top 10 dimensions (indices 548, 527, 401, 590, 511,
1, 505, 287, 442, 124) carry bracket-norm corrected values of
$3.38$--$2.79$, while the median dimension has a corrected value of
$1.42$. These high-bracket-norm dimensions represent the directions
most likely to encode genuine per-nucleotide structural properties
rather than positional artifacts.

\begin{figure}[h]
\centering
\includegraphics[width=0.9\textwidth]{../figures/bracket_norm_per_dim.png}
\caption{Per-dimension bracket-norm corrected values for the 640
RNA-FM embedding dimensions on the HTT CAG region. The distribution
is right-skewed: a minority of dimensions carry disproportionately
high intensive signal. Dimensions above the median (dashed) are
candidates for structural feature extraction.}
\label{fig:bracket_norm}
\end{figure}

\subsection{Grassmannian transportability collapses for any repeat-count change}

The wild-type representation subspace ($r = 21$) is dominated by its
first principal component, which explains 95.7\% of variance. This
near-one-dimensional structure means the $k=10$ PCA subspace is
effectively rank-1: nine of the ten basis vectors capture less than
0.2\% of variance each.

Consequently, any deviation from the wild-type repeat count produces
a near-orthogonal subspace. Geodesic distances are uniformly large:
$d_g = 3.70$ ($r = 15$), $3.71$ ($r = 27$), $3.64$ ($r = 36$),
$3.70$ ($r = 40$), $3.68$ ($r = 50$), $3.79$ ($r = 60$), $3.53$
($r = 80$). Maximum principal angles range from $85.7^\circ$ to $89.9^\circ$
across all variants. The control comparison ($r = 21$ vs.\ itself)
yields $d_g = 0.14$ with a maximum angle of $6.8^\circ$, confirming that
the large distances reflect genuine representation shifts rather
than numerical noise.

\begin{figure}[h]
\centering
\includegraphics[width=0.9\textwidth]{../figures/transportability.png}
\caption{(a)~Geodesic distance from the WT ($r=21$) PCA subspace
to each variant. The distance is uniformly $\sim 3.7$ for all non-WT
repeat counts, with no gradual degradation or phase transition at the
disease threshold ($r \approx 36$). (b)~Principal angle heatmap:
nearly all angles exceed $60^\circ$, indicating near-complete subspace
rotation.}
\label{fig:transport}
\end{figure}

This result has direct implications for RNA-targeting drug design.
Structural features of the HTT CAG region identified from RNA-FM
embeddings at one repeat length do not transfer to any other repeat
length: the representation geometry undergoes a categorical shift,
not a gradual drift. ASO design pipelines that rely on embedding
similarity between wild-type and disease-length variants are
confounded by this length-dependent subspace rotation.

\subsection{Extreme representational compression in later layers}

RNA-FM's 13 layers (embedding + 12 transformer) progressively compress
the representation from 51 effective dimensions (layers~0--2) to a
single effective dimension (layers~9--12). The compression trajectory
is monotonic: $51 \to 51 \to 51 \to 44 \to 30 \to 20 \to 11 \to 8
\to 4 \to 1 \to 1 \to 1 \to 1$. By layer~11, a single principal
component explains 97.9\% of variance.

\begin{figure}[h]
\centering
\includegraphics[width=0.9\textwidth]{../figures/cross_layer_similarity.png}
\caption{Cross-layer subspace similarity (mean cosine of principal
angles between 10-dimensional PCA subspaces). Early layers (0--2) form
a coherent block, as do the late layers (9--12). The transition region
(layers~3--8) shows progressive divergence from the early block and
convergence toward the late block.}
\label{fig:layers}
\end{figure}

This compression is far more extreme than observed in GPT-2, where
the effective dimensionality remains $>10$ through the final layer.
The collapse to rank-1 in RNA-FM suggests that the model's pre-training
objective (masked nucleotide prediction) forces late-layer
representations to converge to a single discriminative direction,
losing the multi-dimensional structural information present in
middle layers. For downstream structural analysis, middle layers
(4--8) may carry richer representations than the conventionally used
final layer.

\subsection{Weight matrices admit progressive low-rank factorization}

SVD of the six projection matrices per layer reveals systematic
structure. The query and key matrices show progressive rank
compression: $W_Q$ effective rank decreases from 265 (layer~0) to 86
(layer~11); $W_K$ follows a parallel trajectory from 254 to 89. This
$3\times$ compression in the QK circuit mirrors the pattern observed
in GPT-2 and suggests that later attention layers attend over a
progressively smaller subspace of the input.

The value matrix $W_V$ maintains high effective rank ($\sim$370)
across all layers, consistent with its role in passing full
representational content through the OV circuit. $W_O$ effective rank
remains moderate ($\sim$300--360). MLP projections $W_{\mathrm{in}}$
and $W_{\mathrm{out}}$ show the highest effective ranks (447--560),
operating over a larger fraction of the full 640-dimensional space.

\begin{figure}[h]
\centering
\includegraphics[width=0.9\textwidth]{../figures/weight_factorization.png}
\caption{(a)~Effective rank (95\% cumulative variance) by layer and
projection type. $W_Q$ and $W_K$ compress $3\times$ from early to late
layers; $W_V$ remains high-rank throughout.
(b)~Spectral decay ratio $\sigma_1/\sigma_{10}$: higher values indicate
stronger concentration of variance in the leading singular directions.}
\label{fig:factorize}
\end{figure}

Cross-projection subspace overlap reveals shared structure within the
QK circuit. The mean cosine of principal angles between the top-20
right singular vectors of $W_Q$ and $W_K$ reaches 0.81 in layer~3 and
0.77 in layer~4, indicating that these projections read from a largely
shared input subspace. By contrast, $W_Q$--$W_O$ overlap is
consistently low ($\sim$0.15), and $W_Q$--$W_V$ overlap is moderate
in layer~0 (0.48) but drops below 0.18 in layers~1--11. This pattern
supports a shared factor bank decomposition: $W_Q$ and $W_K$ can be
decomposed into a common set of input-space factors with
projection-specific selector matrices, as demonstrated for GPT-2 in
prior work.

\subsection{CAG repeat region occupies a distinct representational subspace}

PCA of final-layer embeddings for a 263-nucleotide window (100~nt
flanking + 63~nt CAG repeat + 100~nt flanking) reveals that PC1
explains 98.6\% of variance. Despite this extreme compression, the
CAG repeat positions and flanking positions separate cleanly along
PC1, occupying distinct regions of the 1-dimensional subspace.

\begin{figure}[h]
\centering
\includegraphics[width=0.9\textwidth]{../figures/position_analysis.png}
\caption{(a)~PCA of final-layer RNA-FM embeddings for the HTT CAG
region plus flanking sequence. CAG repeat positions (red) separate
from flanking positions (blue) along PC1 (98.6\% variance).
(b)~Position-wise embedding norm is constant at $28.667$ across all
positions, confirming that LayerNorm equalizes magnitudes and
structural information resides entirely in direction.}
\label{fig:positions}
\end{figure}

The per-position embedding norm is identical for CAG and flanking
positions ($28.667$ in both cases), confirming that the geometric
distinction is directional rather than magnitude-based. This is
consistent with the LayerNorm architecture and underscores the
importance of subspace-based analyses over norm-based ones for
RNA foundation model embeddings.


%% ────────────────────────────────────────────
\section{Discussion}
\label{sec:discussion}

\subsection{Implications for RNA-targeting drug design}

The transportability analysis reveals a categorical failure mode: RNA-FM
representations of the HTT CAG region undergo near-complete subspace
rotation ($d_g \approx 3.7$, principal angles $> 85^\circ$) for any change
in repeat count, with no gradual degradation or phase transition at the
disease threshold. This is a stronger negative result than expected.
Rather than a smooth curve from which one might interpolate, the
geodesic distance is effectively binary: same repeat count
($d_g = 0.14$) or different repeat count ($d_g \approx 3.7$).

This finding has two practical implications. First, ASO screening
pipelines that use RNA-FM embedding similarity between wild-type and
mutant sequences to predict mechanism transfer are confounded: the
observed distances reflect repeat-count-dependent subspace rotation
rather than genuine structural divergence. Second, the bracket-norm
correction identifies which embedding dimensions encode
intensive (per-nucleotide) structural properties that may be more
robust to repeat-count variation. Drug design should target these
intensive dimensions rather than relying on aggregate subspace
geometry.

\subsection{Weight factorization and mechanistic interpretability}

The SVD analysis confirms that RNA-FM's weight matrices have the
low-rank structure required for shared factor bank decomposition. The
$W_Q$ effective rank of 86 in layer~11 (out of 640 dimensions)
implies that a factor bank of $\sim$100--150 learned directions could
capture 95\% of the query-key computation. The high $W_Q$--$W_K$
subspace overlap (0.81 in layer~3) means these factors would be
genuinely shared across projections, yielding an interpretable
decomposition rather than a redundant one.

A factorized RNA-FM would enable direct inspection of which learned
structural motifs (factors) each attention head recruits. The
progressive QK compression from early to late layers suggests a
hierarchy: early layers attend broadly over the full input subspace
($\sim$260 effective dimensions), while late layers attend over a
narrow, task-specialized subspace ($\sim$90 effective dimensions).
Identifying the specific directions that survive this compression
would reveal what RNA-FM considers the most discriminative structural
features --- directly relevant to understanding which aspects of RNA
secondary structure drive the model's predictions.

\subsection{Limitations}

Several limitations qualify these results. RNA-FM was pre-trained
on non-coding RNA, not mRNA; its representations of protein-coding
sequences like HTT may not capture coding-specific structural features.
The PCA-based subspace extraction assumes linearity in the
representation geometry; nonlinear structure would require manifold
learning approaches. The synthetic repeat variants do not capture
somatic mosaicism or epigenetic modifications present in patient
tissue. The weight SVD characterizes the model's \emph{capacity}
for low-rank structure, but the full distillation-based factorization
(fitting a shared factor bank and per-projection selectors) requires
GPU training that we defer to subsequent work.

\subsection{Future directions}

Three extensions follow directly. First, cross-view validation using
thermodynamic structure prediction (EternaFold, RNAfold) as an
independent view, assessing whether the Grassmannian subspace identified
by RNA-FM aligns with physics-based structural ensembles. Second, full
distillation-based factorization of RNA-FM to extract interpretable
factors, following the pipeline established for
GPT-2~\citep{tower2026fac}. Third, application
to additional RNA therapeutic targets (SARS-CoV-2 frameshifting element,
DMD exon-skipping targets, SMA survival motor neuron) to test whether
the bracket-norm and transportability metrics generalize across RNA
drug targets.


%% ────────────────────────────────────────────
\section{Conclusion}
\label{sec:conclusion}

We applied three geometric audits---bracket-norm correction,
Grassmannian transportability, and weight factorization---to RNA-FM
representations of the HTT CAG repeat region.
The central finding is a categorical failure of naive transportability:
RNA-FM embeddings undergo near-complete subspace rotation for any
change in CAG repeat count, with no gradual degradation or clinically
relevant threshold effect. Weight factorization confirms that the
model's internal structure admits low-rank decomposition ($W_Q$
effective rank 86 in layer~11, $W_Q$--$W_K$ overlap 0.81 in
layer~3), providing a path toward mechanistic interpretability of
RNA foundation models. These tools together establish a geometric
validity check for RNA structural claims derived from foundation
model embeddings, and quantify the transportability barrier that must
be overcome for embedding-guided RNA-targeting drug design.

All code, data, and analysis results are available at
\url{https://github.com/ElliottTower/causal-rna} under the MIT license.


%% ────────────────────────────────────────────
\bibliographystyle{plainnat}
% \bibliography{refs}

\begin{thebibliography}{99}

\bibitem[Tower(2026a)]{tower2026bn}
Tower, E.
\newblock Bracket norm: Distinguishing intensive from extensive properties in high-dimensional biomedical feature spaces.
\newblock \emph{Zenodo}, 2026.

\bibitem[Tower(2026b)]{tower2026grass}
Tower, E.
\newblock Grassmannian cohort transfer: Quantifying subspace transportability across biomedical conditions.
\newblock \emph{Zenodo}, 2026.

\bibitem[Tower(2026c)]{tower2026fac}
Tower, E.
\newblock Factorization circuits: Reading transformer circuits from weight decompositions.
\newblock \emph{Zenodo}, 2026.

\bibitem[Chen et al.(2022)]{rnafm}
Chen, J., Hu, Z., Sun, S., Tan, Q., Wang, Y., Yu, Q., Zong, L., Hong, L., Xiao, J., King, I., et al.
\newblock Interpretable RNA foundation model from unannotated data for highly accurate RNA structure and function predictions.
\newblock \emph{arXiv:2204.00300}, 2022.

\bibitem[de~Mezer et al.(2011)]{demezer2011}
de~Mezer, M., Wojciechowska, M., Napierala, M., Sobczak, K., and Krzyzosiak, W.~J.
\newblock Mutant {CAG} repeats of {Huntingtin} transcript fold into hairpins, form nuclear foci and are targets for {RNA} interference.
\newblock \emph{Nucleic Acids Research}, 39(9):3852--3863, 2011.

\bibitem[Devlin et al.(2019)]{devlin2019bert}
Devlin, J., Chang, M.-W., Lee, K., and Toutanova, K.
\newblock {BERT}: Pre-training of deep bidirectional transformers for language understanding.
\newblock In \emph{NAACL-HLT}, 2019.

\bibitem[Edelman et al.(1998)]{edelman1998}
Edelman, A., Arias, T.~A., and Smith, S.~T.
\newblock The geometry of algorithms with orthogonality constraints.
\newblock \emph{SIAM Journal on Matrix Analysis and Applications}, 20(2):303--353, 1998.

\bibitem[MacDonald et al.(1993)]{macdonald1993}
MacDonald, M.~E., Ambrose, C.~M., Duyao, M.~P., et al.
\newblock A novel gene containing a trinucleotide repeat that is expanded and unstable on {Huntington's} disease chromosomes.
\newblock \emph{Cell}, 72(6):971--983, 1993.

\bibitem[Rubinsztein et al.(1996)]{rubinsztein1996}
Rubinsztein, D.~C., Leggo, J., Coles, R., et al.
\newblock Phenotypic characterization of individuals with 30--40 {CAG} repeats in the {Huntington} disease ({HD}) gene reveals {HD} cases with 36 repeats and apparently normal elderly individuals with 36--39 repeats.
\newblock \emph{American Journal of Human Genetics}, 59(1):16, 1996.

\bibitem[Sobczak et al.(2003)]{sobczak2003}
Sobczak, K., de~Mezer, M., Michlewski, G., Krol, J., and Krzyzosiak, W.~J.
\newblock {RNA} structure of trinucleotide repeats associated with human neurological diseases.
\newblock \emph{Nucleic Acids Research}, 31(19):5469--5482, 2003.

\bibitem[Tabrizi et al.(2019)]{tabrizi2019}
Tabrizi, S.~J., Leavitt, B.~R., Landwehrmeyer, G.~B., et al.
\newblock Targeting {Huntingtin} expression in patients with {Huntington's} disease.
\newblock \emph{New England Journal of Medicine}, 380(24):2307--2316, 2019.

\bibitem[Tabrizi et al.(2022)]{tabrizi2022}
Tabrizi, S.~J., Estevez-Fraga, C., van Roon-Mom, W.~M.~C., et al.
\newblock Potential disease-modifying therapies for {Huntington's} disease: lessons learned and future opportunities.
\newblock \emph{The Lancet Neurology}, 21(7):645--658, 2022.

\bibitem[Wolf et al.(2020)]{wolf2020transformers}
Wolf, T., Debut, L., Sanh, V., et al.
\newblock Transformers: State-of-the-art natural language processing.
\newblock In \emph{EMNLP: System Demonstrations}, 2020.

\end{thebibliography}

\end{document}
